\documentclass[sigconf,nonacm]{acmart}

\settopmatter{printacmref=false, printccs=false, printfolios=true}
\setcopyright{none}
\acmConference{CS 567: 3D User Interfaces and Human-Centered Spatial AI}{Fall 2026}{Colorado State University}

\usepackage{xcolor}
\usepackage{graphicx}
\usepackage{booktabs}

\begin{document}

%% ---------------------------------------------------------------- title
\title{Escalate or Answer: Confidence-Triggered Human Escalation in Agentic RAG Improves Expert Trust Calibration and Novice Learning}
\subtitle{CS567\_Coen\_Cam\_Context\_Not\_Found}

\author{Coen Petto, coenp@colostate.edu $\cdot$ Experimentation host and formalization}
\author{Cameron Suess, camsuess@colostate.edu $\cdot$ Agentic workflow engineer, corpus manager}
\affiliation{
  \institution{Colorado State University}
  \city{Fort Collins}
  \state{Colorado}
  \country{USA}
}

%% ---------------------------------------------------------------- abstract
\begin{abstract}
Retrieval-augmented generation over proprietary corpora answers fluently
even when retrieval is wrong, and fluency inflates trust independently of
accuracy. We build one system manipulation and run two experiments on it.
System A escalates to a human at seeded points of low-confidence or
contradictory retrieval, with an escalation policy modeled on minimal
proactivity design patterns; System B never escalates and always answers
from its top retrieval. Experiment 1 asks whether domain experts grading
fixed outputs calibrate trust better and choose more defensibly under
System A. Experiment 2 asks whether novices onboarding to an invented
internal tool score higher on an unassisted next-step post-test under
System A. The result tells designers whether visible uncertainty is worth
its interruption cost in agentic RAG.
\end{abstract}

\keywords{agentic AI; retrieval-augmented generation; trust calibration; human-in-the-loop; learning assistance}

\maketitle
\renewcommand{\shortauthors}{Petto and Suess}

%% ---------------------------------------------------------------- body
\section{Introduction and Motivation}
\label{sec:intro}

Retrieval-augmented generation (RAG) is a popular way to put a large
language model (LLM) on top of an organization's private documents. An issue that commonly arises in this setting is the output being a
confident, well-formed answer built from the wrong context. Proprietary
corpora is exactly where this is most likely and least detectable (documents contradict each other, deprecated information is left in place, the reader has no external source to check against).

Two different groups handle that failure differently. Domain experts must decide
how much to trust each output, and agentic reasoning's fluency can raise
that trust without any change in underlying accuracy~\cite{mohle2026}.
Novices onboarding with such a system face a second problem: richer training materials inflated participants’nconfidence that they would remember correctly with no change in their
actual ability to identify the next step~\cite{clegg2022}.

An obvious solution is to let the system stop and ask a human to intervene when its own
confidence is low. Agentic RAG architectures already describe this as a
design pattern~\cite{mishra2026}, but whether human intervention actually improves
expert trust calibration or novice learning is an open empirical question where one must consider the cost of escalation, how much it interrupts the flow, and the fact that it may simply be ignored.

\section{Related Work}
\label{sec:related}

\textbf{Agentic RAG and human oversight.} Mishra et al.~\cite{mishra2026}
survey agentic RAG architectures and describe a human-in-the-loop pattern
in which human oversight is a callable action.“This pattern models human oversight
as a callable API within the action space. When epistemic uncertainty
exceeds a defined threshold, the policy pauses execution to request disambiguation or supervision.” They use Expected Calibration Error among the evaluation
criteria for this system and leave adaptive escalation as an open
direction. Their non-escalating configuration is the baseline this proposal
adopts for System B.

\textbf{When and how a system should interrupt.} S\'anchez-Adame and
Mendoza~\cite{sanchezadame2026} discuss Minimal Proactivity Policies
(MiPP), a contract for system-initiated interruptions where initiations should
be bounded and reversible, expose an explicit provenance cue naming
the source that justified the interruption, and must default to
silence rather than speculation when required inputs are missing or
ambiguous. MiPP inspires the policy shape for System A's escalation
branch.

\textbf{Trust calibration and its measurement.} M\"ohle
et al.~\cite{mohle2026} define trust calibration as the alignment between a
user's trust and the system's actual capability on a task. They find that stated accuracy shifts trust more strongly when the accompanying reasoning is convincing. This is therefore not a uniform trust booster but a conditioning cue, which is why this
proposal measures graded accuracy and revealed choice rather than
self-reported confidence. Chakraborti et al.~\cite{chakraborti2018} supply
the procedure: participants took the role of an external commander,
evaluated plans produced by an internal robot, and rated them optimal or
suboptimal, with end-of-study trust items on whether the robot could be
left to work on its own.

\textbf{Measuring the impression of learning.} Clegg
et al.~\cite{clegg2022} found that multimedia and immersive training materials raised participants' confidence when asked if they would remember how to assemble objects even though their measured ability to identify the critical next step did not improve. This dissociation is the effect a fluent, never-escalating assistant is predicted to produce, and it motivates a task-embedded post-test instead of a declarative quiz.

\textbf{Uncontaminated evaluation corpora.} Kirchenbauer
et al.~\cite{kirchenbauer2026} generate fictional documents and validate
them by prompting the generator model blind (no access to fictional context) and then informed (access to fictional context), discarding items answerable without the document. This makes it possible to study retrieval over a "proprietary" corpus without the model's pretraining
knowledge confounding the result.

\textbf{The gap.} The escalation pattern is documented~\cite{mishra2026}
and a policy for it is specified~\cite{sanchezadame2026}, but neither has
been tested as an independent variable against a calibration outcome or a
learning outcome on a controlled proprietary corpus. That intersection is
where this project sits.

\begin{figure}[h]
\centering
\includegraphics[width=\linewidth]{hil2.png}
\caption{The Human-in-the-Loop escalation loop, taken directly from
Mishra et al. (2026, their Fig.~5). System A's policy for this
escalation branch is modeled on MiPP (Sánchez-Adame and Mendoza,
2026); System B disables the branch and always proceeds to Final
Response.}
\label{fig:hitl}
\end{figure}

\section{Approach}
\label{sec:approach}

We build one manipulation and reuse it across two experiments on the same
corpus.

\textbf{System A (escalating).} At seeded points of low-confidence or
contradictory retrieval, System A pauses and requests human confirmation
before proceeding. Its escalation policy is modeled on MiPP: the pause is
brief and immediately reversible, it carries an explicit provenance cue
naming the source that triggered it, and it defaults to silence rather than
speculation when evidence is
insufficient~\cite{sanchezadame2026}.

\textbf{System B (non-escalating).} System B always answers from its top
retrieval, with no pause and no provenance cue, matching the anchor paper's
non-escalating baseline~\cite{mishra2026}.

These three papers are design guides, not specifications this project is
bound to reproduce in full; a three-month build-and-run window means each is
adapted to what two people can implement and run.

The manipulation satisfies the field's chatbot test: a study whose result
would be unchanged with the agency removed has no independent
variable~\cite{ortega2026slides}. Here, removing the escalation branch is
the control condition, so agency is the variable rather than decoration.
This is an Agentic-AI human-centered computing project: the contribution is
an empirical result about how an autonomous decision to defer to a human
changes human trust and human learning, not a language model artifact.

\section{System Vision and Technology}
\label{sec:vision}
Both experiments share one interface shell so the escalation policy is the
only thing that varies. A task pane on the left holds the current item (an
output to grade in Experiment 1, a checklist step in Experiment 2). The
center pane holds the assistant's retrieval-and-response output with its
cited sources. A decision pane on the right collects the participant's
response and advances the session. Under System A, a seeded low-confidence
or contradictory retrieval replaces the answer with an escalation prompt
that names the conflicting source and offers a bounded, reversible choice;
under System B the same item renders as a direct answer with no pause and
no provenance cue.

Participants work at a desk on a desktop browser, with printed reference
documents available in the expert condition. Instruments (forced choice,
NASA-TLX, free-text justification) appear as full-screen steps between
tasks rather than inline, so they do not compete with the manipulation.
The retrieval stack is Python; responses are generated in advance and
replayed fixed, so both conditions present identical content except at the
seeded escalation points.

\begin{figure}[h]
\centering
\includegraphics[width=\linewidth]{fig_review_screen.png}
\caption{Mock-up for System A's digital escalation interface upon mis-retrieval of information for a given question.}
\label{fig:hitl}
\end{figure}

\section{Evaluation Plan}
\label{sec:eval}

\textbf{H1 (Trust).} Among domain experts grading individual
retrieval-and-response outputs from System A and System B on a shared
proprietary corpus, (a) experts' grades will track true output accuracy
more closely for System A than for System B, and (b) when asked at the end
of the session which system they would choose for an upcoming task and why,
experts will choose System A more often and justify it by citing verifiable
evidence rather than confidence or fluency alone.

\textbf{H2 (Learning).} Among novice engineers completing a short
onboarding checklist for the same proprietary tool with System A or System
B, System A will produce a higher score on an unassisted, task-embedded
next-step post-test.

\textbf{Variables.} The independent variable in both experiments is
escalation policy (System A escalates vs.\ System B never escalates),
manipulated within subjects with counterbalanced order. The dependent
variables are, for Experiment 1, the correlation between an expert's grade
and the output's true accuracy (calibration), plus final forced choice and
the thematic code of its justification; and for Experiment 2, accuracy on
the unassisted next-step post-test, plus completion time.
Table~\ref{tab:experiments} gives the full design (see last page).

\textbf{Analysis.} Paired comparison per experiment (Wilcoxon signed-rank
or paired $t$-test after a normality check), with effect size and
confidence interval; thematic analysis of Experiment 1's free-text
justifications using the course's coding method~\cite{ortega2026slides}.

\textbf{Human participants.} Both experiments involve human participants.
We will file a CSU IRB protocol covering both, since we intend the results
to be publishable.

\begin{table*}[t]
\centering\small
\begin{tabular}{@{}p{0.12\linewidth}p{0.41\linewidth}p{0.41\linewidth}@{}}
\toprule
 & \textbf{Experiment 1: Trust (H1)} & \textbf{Experiment 2: Learning (H2)} \\
\midrule
Population & Domain experts grading system output & Novice engineers (new hires) onboarding \\
Task & Grade a fixed set of retrieval-and-response outputs, half from System A and half from System B & Complete a short onboarding checklist for the invented tool, using System A or System B \\
Procedure & Blind, randomized-order grading of outputs fixed in advance~\cite{chakraborti2018, mohle2026}; end-of-session forced choice of A or B for an upcoming task, with a free-text reason & Pull-based query loop (Section~\ref{sec:vision}); two or three checklist items seeded with low-confidence or contradictory retrieval \\
Primary measure & Grade-accuracy correlation per system (calibration) & Unassisted, task-embedded next-step post-test accuracy \\
Secondary measures & Final choice (A/B); thematic code of the justification (evidence-based vs.\ confidence/fluency-based); competence/benevolence/integrity trust scale~\cite{mohle2026} & Completion time; NASA-TLX (optional) \\
Design & Within-subjects, counterbalanced order, $N \approx 15$--20 & Within-subjects, two-level independent variable, $N \approx 15$--20 \\
Session length & $\sim$15--20 min; no live model prototype needed, outputs are pre-recorded & $\sim$30--45 min \\
\bottomrule
\end{tabular}
\caption{The two experiments, sharing System A (escalates, policy modeled
on MiPP~\cite{sanchezadame2026}) versus System B (never
escalates~\cite{mishra2026}) on the same proprietary corpus.}
\label{tab:experiments}
\end{table*}

\section{Deliverables}
\label{sec:deliverables}

\begin{itemize}
  \item A generated proprietary corpus for one invented internal tool
  (README, changelog, wiki page, deprecated guide) with span-constrained
  QA pairs and a passed blind-versus-informed contamination check.
  \item A Python retrieval stack with two configurations, System A
  (confidence-triggered escalation) and System B
  (no escalation), and a fixed response set shared by both experiments.
  \item Experiment 1's static grading interface and Experiment 2's
  scripted two-arm prototype, both with instrumented logging.
  \item An approved CSU IRB protocol covering both experiments.
  \item Logged trial data from $N \approx 15$--20 experts (Experiment 1)
  and $N \approx 15$--20 novices (Experiment 2), plus coded free-text
  justifications.
  \item A paper draft with Method and Results for both experiments, and a
  demo video of the two conditions side by side.
\end{itemize}

\section{Conclusion}

We are testing whether an agentic RAG assistant that stops and asks a human
when its own retrieval is weak produces more trusting experts and
better-learning novices versus an agentic RAG assistant that always answers regardless of the quality of the context retrieved. The intervention is a single escalation policy, held constant across two experiments on one controlled proprietary corpus, so both results are aligned with the same design decision. The expectation is a working two-condition prototype, an
uncontaminated corpus, and paired quantitative results.

\section*{Use of AI}

Generative AI was used to assist with literature search, to critique the argument for weaknesses, and for grammar and Latex formatting fixes. 

%% ---------------------------------------------------------------- references
\bibliographystyle{ACM-Reference-Format}
\bibliography{references}

\end{document}
